\tzpanchor{0506}
\tzpentryheader{0506}{PURCELL FACTOR \& pi-MODULATED CAVITY MODES}

\begingroup\small
\begingroup\scriptsize
\setlength{\tabcolsep}{4pt}
\renewcommand{\arraystretch}{0.92}
\noindent\begin{tabularx}{\linewidth}{@{}p{0.24\linewidth}p{0.36\linewidth}X@{}}
\textbf{UUID} & \multicolumn{2}{l@{}}{\tzpcode{ba4f3275-56da-52f9-ae6f-e93929852ffe}}\\
\textbf{PiID} & \tzpcode{3362440656643} & \textbf{TZPi}\quad \tzpcode{0}\quad \textbf{PiPE}\quad \tzpcode{513}\\
\textbf{RTEID} & \textbf{Poloidal $\theta$}\quad \tzpcode{1945.4943 rad} & \textbf{Toroidal $\phi$}\quad \tzpcode{03.2056 rad}\\
\textbf{TZPID} & \tzpcode{ID0506} & \textbf{Node Dependencies}\quad \tzpidref{0505}\\
\textbf{Registry} & \tzpcode{registry\_id\_0506} & \textbf{Scale}\quad transfer\\
\textbf{LOC Classification} & QA641 manifolds and topology & QC174.17 mathematical physics\\
\textbf{arXiv Classification} & Primary: math-ph & Cross-list: gr-qc, math.DG\\
\textbf{Primary Support} & Purcell Factor & Pi-Mode Constraint\\
\textbf{Closure Support} & Cavity Closure & \\
\end{tabularx}\par
\endgroup
\endgroup

\tzpsection{Canonical Statement}
F\_p = frac\{3\}\{4pi\textasciicircum{}2\} left(frac\{lambda\}\{n\}right)\textasciicircum{}3 frac\{Q\}\{V\} otimes text\{kissingNumber\}

\tzpsection{Canonical Equation}
\[
F_{\mathrm{p}} = \frac{3}{4\pi^2}\left(\frac{\lambda}{n}\right)^3 \frac{Q}{V}
\]
\[
phi_m = (2m+1)\pi
\]

\begin{minipage}[t]{0.49\linewidth}
\tzpsection{Canonical Symbol Key}
\begingroup
\small
\raggedright
\textbf{Source equation symbols.}\par
\begin{itemize}[leftmargin=1.35em,itemsep=1pt,topsep=1pt,parsep=0pt]
    \item $F_{\mathrm{p}}$: source equation symbol.
    \item $\pi$: source equation symbol.
    \item $\lambda$: cosmological or lemniscatic branch parameter in the entry relation.
    \item $n$: integer mode index.
    \item $Q$: source equation symbol.
    \item $V$: auxiliary target state or comparison variable.
    \item $\mathcal{M}_{\pi}$: source equation symbol.
\end{itemize}
\endgroup
\end{minipage}\hfill
\begin{minipage}[t]{0.47\linewidth}
\tzpsection{Wolfram Form}
\begingroup
\scriptsize
\raggedright
\textbf{Source Wolfram model.}\par
\begin{Verbatim}[fontsize=\tiny,breaklines=true,breakanywhere=true]
(* TZPID ID0506 generated source Wolfram model *)
ClearAll[K0506, Cr0506, T, R, A, W, I, N];
eq0506$1 = HoldForm[FP == (3/(4*Pi^2))*(lambda/n)^3*(Q/V)];
eq0506$2 = HoldForm[phiM == (2*m + 1)*Pi];
eq0506$3 = HoldForm[mathcal[A]*_*adm*MPi == 1];

K0506 = HoldForm[{eq0506$1, eq0506$2, eq0506$3}];

N = "normalized TRAWIN closure object for PURCELL FACTOR \& pi-MODULATED CAVITY MODES";
I = "Cavity Closure integration/comparison layer";
W = "Pi-Mode Constraint wave or operator carrier";
A = "Purcell Factor admissible action/amplitude condition";
R = "Pi-Mode Constraint rotation/curl preservation layer";
T = "Purcell Factor local source condition";

Cr0506[expr_] := HoldForm[N[I[W[A[R[T[expr]]]]]]];

Result0506 = Cr0506[K0506];
\end{Verbatim}
\endgroup
\end{minipage}

\par\vspace{0.45\baselineskip}
\noindent\begin{minipage}[t]{\linewidth}
\begingroup
\small
\raggedright
\textbf{TRAWIN Operator Key.}\par
\noindent\textbf{Time/localization operator:} $\mathsf{T}\!\left[\mathrm{local}\right]$ Purcell Factor local source condition.\par
\noindent\textbf{Rotation/curl operator:} $\mathsf{R}\!\left[\nabla\times\right]$ Pi-Mode Constraint rotation/curl preservation layer.\par
\noindent\textbf{Action/amplitude operator:} $\mathsf{A}\!\left[\mathrm{admissible}\right]$ Purcell Factor admissible action/amplitude condition.\par
\noindent\textbf{Wave/d'Alembertian operator:} $\mathsf{W}\!\left[\Box\right]$ Pi-Mode Constraint wave or operator carrier.\par
\noindent\textbf{Integration operator:} $\mathsf{I}\!\left[\int\right]$ Cavity Closure integration/comparison layer.\par
\noindent\textbf{Normalization operator:} $\mathsf{N}\!\left[\mathrm{normalized}\right]$ normalized TRAWIN closure object for PURCELL FACTOR \textbackslash{}\& pi-MODULATED CAVITY MODES.\par
\endgroup
\end{minipage}

\clearpage
\noindent\textbf{[ID 0506] PURCELL FACTOR \& pi-MODULATED CAVITY MODES}\par
\vspace{0.35em}
\tzpsection{Derivation Layer}
\paragraph{Primary Equation.}
\begin{equation}
u_T := \int_{\Gamma}\sqrt{|g|}\,\mathrm{d}\ell .
\end{equation}

\paragraph{Primary Derivation.}
Let $\Gamma$ be a fundamental loop encircling the Trawin Zero Point exactly once.  
The infinitesimal line element induced by the metric is:

\begin{equation}
\mathrm{d}s = \sqrt{|g|}\,\mathrm{d}\ell .
\end{equation}

Integrating this metric-weighted length around the loop gives the primitive unit associated with one complete encirclement of the TZP:

\begin{equation}
u_T = \int_{\Gamma}\mathrm{d}s .
\end{equation}

Substituting the line element yields:

\begin{equation}
\boxed{
u_T := \int_{\Gamma}\sqrt{|g|}\,\mathrm{d}\ell .
}
\end{equation}

\paragraph{Interpretation.}
The Trawin base unit $u_T$ is the fundamental geometric unit obtained by completing one metric-weighted loop around the TZP.

\par\vspace{0.35em}
\noindent\textbf{Derivable Consequences.}\quad
\begingroup
\footnotesize
\raggedright
The canonical equation anchors PURCELL FACTOR \textbackslash{}\& pi-MODULATED CAVITY MODES by connecting the source condition to the derivation layer and proof certificate.
\par\endgroup

\tzpsection{Equation Support}
\noindent\textbf{1. Purcell Factor}\par
\[
F_{\mathrm{p}} = \frac{3}{4\pi^2}\left(\frac{\lambda}{n}\right)^3 \frac{Q}{V}
\]
\noindent This fixes the cavity-enhancement factor carried by the resonant mode.\par
\noindent\textbf{2. Pi-Mode Constraint}\par
\[
phi_m = (2m+1)\pi
\]
\noindent This imposes the odd-\textbackslash{}textbackslash\{\}pi phase ladder of the modulated cavity branch.\par
\noindent\textbf{3. Cavity Closure}\par
\[
mathcal{A}_{\mathrm{adm}}\!\left(\mathcal{M}_{\pi}\right)=1
\]
\noindent This closes the entry by requiring admissibility of the resulting pi-mode sector.\par

\clearpage
\noindent\textbf{[ID 0506] PURCELL FACTOR \& pi-MODULATED CAVITY MODES}\par
\vspace{0.35em}
\tzpsection{Dictionary}
\begingroup\small
\tzpsection{Definition}
PURCELL FACTOR \& pi-MODULATED CAVITY MODES is a registry definition in Quantum-Classical Transition with secondary pressure from Gyromagnetic and Rotating-Frame Physics.

% BEGIN TZP CONTEXTUAL MEANING
\tzpsection{Contextual Meaning}
\begingroup\footnotesize\setlength{\emergencystretch}{3em}\sloppy
\textbf{Formal statement:} F sub p = frac 3 4pi to the power 2 left(frac lambda n right) to the power 3 frac Q V otimes text kissingNumber. \textbf{Contextual meaning:} The entry reads PURCELL FACTOR \& pi-MODULATED CAVITY MODES as a controlled technical headword whose equation layer, proof route, and registry role are interpreted together. \textbf{Derivable consequences:} The entry now carries a clean cavity-enhancement law and an explicit pi-mode constraint instead of a malformed composition token. \textbf{Lemma support:} Purcell Factor; Pi-Mode Constraint; Cavity Closure.\par
\endgroup
% END TZP CONTEXTUAL MEANING

\tzpsection{Technical Interpretation}
The Trawin base unit u sub T is the fundamental geometric unit obtained by completing one metric-weighted loop around the TZP.

\tzpsection{Equation Grounding}
Equation anchors: F sub p = ratio 3 4pi to the power 2 (ratio lambda n ) to the power 3 ratio Q V; phi sub m = (2m+1)pi; and mathcal A sub adm ( M sub pi )=1. Proof route: show that the cavity enhancement and odd-pi phase ladder together determine an admissible modulated mode sector. Formalization route: prove that the pi-modulated branch remains admissible under the stated Purcell enhancement law.

\tzpsection{Interpretive Role}
The entry now carries a clean cavity-enhancement law and an explicit pi-mode constraint instead of a malformed composition token.
\endgroup

\tzpsection{TZP Thesaurus}
\begin{enumerate}[leftmargin=1.6em,itemsep=6pt,topsep=4pt]
\item \textbf{Spontaneous Emission Enhancement}: The physical phenomenon where trapping an atom in a perfectly sized mirrored box forces it to release its energy much faster than it normally would.
\item \textbf{Anti-Symmetric Phase Locking}: The boundary requirement forcing adjacent resonant chambers to operate exactly half a wavelength out of sync with one another.
\item \textbf{High-Q Resonator Tuning}: The extreme engineering precision required to minimize energy leakage and maximize the bouncing, overlapping internal echoes of the cavity.
\end{enumerate}

\tzpsection{Encyclopedia Mathematica}
\begingroup\footnotesize\sloppy
The visual plate below is reserved for the source-specific equation and progression graphic for [ID 0506]. It is included only as a companion to the canonical equation, not as a substitute for the formal text.
\par\endgroup
\ifTZPDarkEdition
\tzpdictionaryvisualband{D:/TZPID/kdp_workflow/build/tikz_figures/ID0506_dictionary_figure_dark.pdf}
\fi

\clearpage
\begingroup\tzpsupportsetup
\noindent\textbf{\fontsize{10.8pt}{12.8pt}\selectfont [ID 0506] Constructive Logic and Proof Certificate}\par
\noindent\textbf{\fontsize{10pt}{12pt}\selectfont PURCELL FACTOR \& pi-MODULATED CAVITY MODES}\par
\vspace{0.45em}
\begingroup
\footnotesize
\setlength{\parskip}{0.22em}
\setlength{\parindent}{0pt}
\noindent\textbf{Curry--Howard--Lambek Interpretation}\par
Under the Curry--Howard--Lambek correspondence, this registry entry is read as a typed proof
object: the stated source condition supplies the proposition, the derivation supplies the
morphism, and the proof certificate records the machine handles needed to check the obligation
in the formal registry.

\vspace{0.45em}
\noindent\textbf{CHL Proof}\par
\textbf{Statement:} f\textbackslash{}\_p = frac\textbackslash{}3\textbackslash{}\textbackslash{}4pi\textasciicircum{}2\textbackslash{} left(frac\textbackslash{})\textasciicircum{}3 frac\textbackslash{}\textbackslash{} otimes text\textbackslash{}

\textbf{Logic Validation:}\\
\textbf{Proposition:} ``PURCELL FACTOR \& \textbackslash{}pi-MODULATED CAVITY MODES is valid as a tensorial/geometric
statement when the stated tensor fields are well formed on the specified domain.''\\
\textbf{Proof:} The source statement gives a metric, curvature, or tensor relation. The CHL reading
validates it by checking that the tensorial objects have compatible domains, indices,
and transformation behavior.

\textbf{VERDICT:} \ensuremath{\checkmark}\ \textbf{TRUE} --- tensorial well-formedness validation

\textbf{Formal Status:} \ensuremath{\checkmark}\ THEORETICAL -- source-backed CHL proof obligation

\textbf{Physical Status:} THEORETICAL -- requires model-specific empirical interpretation
\par\vspace{0.45em}
\noindent{\tiny\textbf{Isabelle/HOL.} \tzpplaincode{physics\_validity\_from\_model, external\_validity\_package, registry\_number\_matches, equation\_count\_matches}}\par
\ifTZPDarkEdition
\tzpproofvisualband{D:/TZPID/kdp_workflow/build/tikz_figures/ID0506_proof_certificate_figure_dark.pdf}
\fi
\endgroup
\endgroup
